%-------------------------------------------------------------------------------
% ResumeOps LaTeX cover letter template
%
% Deliberately matched to resume-template.tex: same fonts, same link colour, same
% header treatment, so the CV and the letter read as one pair of documents.
%
% Content rules (paragraph structure, voice, length) live in
% ../cover-letter-preferences.md. Typography rules live in template-preferences.md.
%
% Everything in <ANGLE BRACKETS> is a placeholder to replace. Delete the date and
% company block if the letter is being pasted into a web form instead of sent.
%
% Build: pdflatex cover-letter.tex   (twice, if links are added)
%-------------------------------------------------------------------------------

\documentclass[letterpaper,10pt]{article}   % 11pt if the letter runs short

\usepackage[empty]{fullpage}
\usepackage[usenames,dvipsnames]{color}
\usepackage{enumitem}
\usepackage{hyperref}
\definecolor{linkblue}{RGB}{0,0,139}
\hypersetup{colorlinks=true, urlcolor=linkblue, linkcolor=linkblue, citecolor=linkblue}
\usepackage{fancyhdr}
\usepackage[english]{babel}                 % spanish, when drafting in Spanish
\usepackage{cmap}          % parsable text layer (must load before fontenc)
\usepackage[T1]{fontenc}
\usepackage[utf8]{inputenc}
\usepackage{lmodern}
\usepackage{helvet}        % header name only
\input{glyphtounicode}

\pagestyle{fancy}
\fancyhf{}
\renewcommand{\headrulewidth}{0pt}
\renewcommand{\footrulewidth}{0pt}
\setlength{\footskip}{12pt}   % silences fancyhdr's \footskip warning

% Roomier than the CV: a letter breathes, a CV does not.
\addtolength{\oddsidemargin}{-0.35in}
\addtolength{\evensidemargin}{-0.35in}
\addtolength{\textwidth}{0.7in}
\addtolength{\topmargin}{-0.55in}
\addtolength{\textheight}{1.1in}

\urlstyle{same}
\raggedright
\setlength{\parindent}{0pt}
\setlength{\parskip}{9pt}     % paragraph spacing instead of indents
\linespread{1.12}

\pdfgentounicode=1            % keep: this is what makes the PDF text extractable

%-------------------------------------------------------------------------------
\begin{document}

%----------HEADER----------
% Same identity block as the CV header, minus the icons — a letter is prose, and
% stray glyphs in the text layer are not worth the visual echo.
\begin{center}
    {\sffamily\bfseries\LARGE <FULL NAME>} \\ \vspace{4pt}
    \small <CITY, COUNTRY> \quad $|$ \quad <+PHONE> \quad $|$ \quad
    \href{mailto:<EMAIL>}{\underline{<EMAIL>}} \quad $|$ \quad
    \href{https://www.linkedin.com/in/<HANDLE>/}{\underline{LinkedIn}}
\end{center}
\vspace{6pt}
\hrule
\vspace{14pt}

%----------DATE + RECIPIENT----------
% Delete this whole block when pasting the letter into an application form.
<DD MONTH YYYY>

\vspace{-4pt}
<COMPANY NAME> \\
<CITY, COUNTRY>

%----------SALUTATION----------
% Name the hiring manager when the JD gives one; otherwise "Dear Hiring Team,".
Dear <HIRING MANAGER / Hiring Team>,

%----------1. HOOK----------
% Role + where seen, then the single strongest reason this candidate fits THIS
% posting. Specific and concrete. Never "I am writing to apply for".
<OPENING: the role by its exact title, and the one fact that makes the fit
obvious — a system owned, a domain, a number.>

%----------2. PROOF----------
% One or two achievements from knowledge/generated/4-achievement-bank.md, told as
% stories rather than re-cut CV bullets: problem, decision, outcome.
<ACHIEVEMENT ONE: the problem, what was decided and built, what changed as a
result — with the metric if one exists in knowledge/.>

<ACHIEVEMENT TWO (optional): a second dimension the JD asks for — scale,
reliability, collaboration, research depth.>

%----------3. FIT----------
% Why this company specifically: their product, domain, stack or public work, and
% how the candidate's direction meets it. If nothing true and specific can be
% said here, ask the user rather than padding it.
<WHY THIS COMPANY: something concrete about them, joined to what the candidate is
deliberately moving towards.>

%----------4. CLOSE----------
% Hard constraints the recruiter must know, then a plain sign-off.
<CLOSE: location / work authorisation / notice period if relevant, and a direct,
unbegging final sentence.>

\vspace{10pt}
Sincerely, \\[6pt]     % "Kind regards," for a warmer or European register
<FULL NAME>

%-------------------------------------------------------------------------------
\end{document}
